\sect{Ćwiczenia 14}

\begin{frame}{Zadanie 1}
	\begin{figure}
		\includegraphics[scale=0.08, alt={}]{./imags/E14/01.png}
		\caption{Zadanie 1. Oblicz widmo sygnału ($a_n$ i $b_n$), wykonaj filtrację w dziedzinie częstotliwości przy pomocy jednego z filtrów z poprzednich ćwiczeń.}
	\end{figure}	
\end{frame}

\begin{frame}{Zadanie 2}
	\begin{figure}
		\includegraphics[scale=0.08, alt={}]{./imags/E14/02.png}
		\caption{Zadanie 2. Oblicz widmo sygnału ($a_n$ i $b_n$), wykonaj filtrację w dziedzinie częstotliwości przy pomocy jednego z filtrów z poprzednich ćwiczeń.}
	\end{figure}	
\end{frame}

\begin{frame}{Zadanie 3}
	\begin{figure}
		\includegraphics[scale=0.08, alt={}]{./imags/E14/03.png}
		\caption{Zadanie 3. Oblicz widmo sygnału ($a_n$ i $b_n$), wykonaj filtrację w dziedzinie częstotliwości przy pomocy jednego z filtrów z poprzednich ćwiczeń.}
	\end{figure}	
\end{frame}

\begin{frame}{Zadanie 4}
	\begin{figure}
		\includegraphics[scale=0.08, alt={}]{./imags/E14/04.png}
		\caption{Zadanie 4. Oblicz widmo sygnału ($a_n$ i $b_n$), wykonaj filtrację w dziedzinie częstotliwości przy pomocy jednego z filtrów z poprzednich ćwiczeń.}
	\end{figure}	
\end{frame}
